\section{Equal width, unequal access: a nonuniform positive split limit}
\label{sec:highorder}
The preceding example concerns missing observations. Even full population
information does not make a fixed structural operation effective. We now
hold the number of children, positivity and sum of outgoing weights, and
weighted incoming mean fixed, but change the relation between the child
weight and the size of the incoming perturbation. This is a restricted
separation, not a new discovery of parity hardness or boundary-weight
escape; see Section~\ref{sec:related}.

Let $d\ge6$ be even, $X$ uniform on $\{-1,1\}^d$, and
$\chi(X)=\prod_{j=1}^d X_j$. For $0<\eta<1/2$, set
\begin{equation}\label{eq:paritylaw}
 \Pr(Y=1\mid X=x)=\tfrac12+\eta\chi(x).
\end{equation}
This law has full support and $I(X;Y)=\log2-h(1/2+\eta)>0$,
independently of $d$. No limiting rare source mass is used when $d$ is fixed.
Fix a smooth activation $\sigma$ near $b$ with
$\sigma''(b)\ne0$ and $\sigma^{(d)}(b)\ne0$, with bounded derivatives
through order $d+2$ there. The experiments use $\sigma=\tanh$ and $b=1/2$.
The parent logit is the constant $\sigma(b)-\sigma(b)=0$.
For a unit vector $v$ put $V=v^\top X$, and define a positive binary split
\begin{equation}\label{eq:asymmetric}
 f_{t,p}(x)=p\sigma(b+tV)+(1-p)\sigma\left(b-\frac{p}{1-p}tV\right)-\sigma(b),
 \quad 0<p<1.
\end{equation}
There are exactly two hidden units. Their outgoing weights sum to one,
their weighted incoming mean is zero, and their outgoing $\ell_1$ norm is
one. For $p\le1/2$, both incoming weights have norm at most $t$.
The intercept stays fixed in this comparison; we do not claim an obstruction
after arbitrary intercept refitting or subsequent parameter training.

Write $R(f)=\E[\log(1+e^{f(X)})-Yf(X)]$ and
\begin{equation}\label{eq:KB}
 K(v)=\frac{\sigma''(b)^2}{32}\E V^4>0,
 \qquad B(v)=\sigma^{(d)}(b)\prod_{j=1}^d v_j.
\end{equation}
A direction with $B(v)>0$ exists by changing the sign of one nonzero
coordinate. Its selection here is an oracle construction.

\begin{theorem}[A constrained weight--radius separation]\label{thm:boundary}
Under the preceding conditions:
\begin{enumerate}[label=(\roman*)]
\item For every fixed $p\in(0,1)$ and unit $v$,
\[
 R(f_{t,p})-\log2
   =K(v)\left(\frac{p}{1-p}\right)^2t^4+o(t^4)>0
\]
for all sufficiently small positive $t$. This positivity is uniform when
$p$ lies in any compact subinterval of $(0,1)$ and $v$ ranges over the unit sphere.
\item Choose $B(v)>0$ and let $p(t)=\kappa t^{d-4}$ for fixed $\kappa>0$.
Then
\begin{equation}\label{eq:jointlimit}
 R(f_{t,p(t)})-\log2
 =\{K(v)\kappa^2-\eta B(v)\kappa\}t^{2d-4}
    +o(t^{2d-4}).
\end{equation}
In particular, the choice
\[\kappa=\frac{\eta B(v)}{2K(v)}\]
gives strict descent with leading coefficient
\[ -\frac{\eta^2B(v)^2}{4K(v)}.\]
\item At the constant parent, the residual splitting matrix is zero.
More generally the derivatives in incoming weights of the label-dependent
part of the population loss agree with the independent-label law through
every order strictly below $d$.
\end{enumerate}
\end{theorem}

\begin{proof}
For every integer $j<d$, $\E[\chi(X)V^j]=0$. Indeed expand $V^j$
into monomials: at least one coordinate has even exponent (possibly zero),
and multiplication by $\chi$ makes its expectation zero. At degree $d$,
the only surviving monomials use each coordinate once, so
\begin{equation}\label{eq:fouriermoment}
 \E[\chi V^d]=d!\prod_i v_i.
\end{equation}
Odd-degree moments with $\chi$ vanish here because $d$ is even and the law
is invariant under $X\mapsto-X$.
The exact loss identity is
\begin{equation}\label{eq:parityrisk}
 R(f)-\log2=\E\log\cosh(f/2)-\eta\E[\chi f].
\end{equation}
For fixed $p$, Taylor expansion of \eqref{eq:asymmetric} cancels the constant
and linear terms and gives, uniformly on the finite input set,
\[
 f_{t,p}=\frac{\sigma''(b)}2\frac{p}{1-p}t^2 V^2+O(t^3).
\]
Since $\log\cosh(z/2)=z^2/8+O(z^4)$, its convexity penalty has leading
term $K(v)[p/(1-p)]^2t^4$. Equations~\eqref{eq:fouriermoment}
and the Taylor expansion through order $d$ show that $\E[\chi f_{t,p}]=O(t^d)$.
Thus the signal is $o(t^4)$. Uniformity on compact $p$ intervals follows
from bounded derivatives and $|V|\le\sqrt d$; $\E V^4\ge(\E V^2)^2=1$
ensures a uniform positive leading coefficient.

For the joint limit it is important to retain the factor $p$ in the
remainders. Put $q=p/(1-p)$. For $p\le1/2$, the two Taylor remainders of
order three are bounded by $C[p+(1-p)q^3]t^3=O(pt^3)$.
Consequently, as both $p,t\to0$,
\[
 f_{t,p}=\tfrac12\sigma''(b)\,p t^2 V^2\{1+O(p)\}+O(pt^3),
\]
where the braces apply to the displayed leading term, not to a pointwise
relative error at $V=0$. It follows that
\begin{equation}\label{eq:penaltyjoint}
 \E\log\cosh(f_{t,p}/2)=K(v)p^2t^4\{1+O(p)+O(t)\}.
\end{equation}
Expanding through order $d+1$, all lower-degree terms and the degree-$d+1$
term disappear after multiplying by $\chi$. The remaining Taylor errors
are $O(pt^{d+2})$. Hence
\begin{align*}
 \E[\chi f_{t,p}]
 &=B(v)t^d\left\{p+(1-p)\left(\frac p{1-p}\right)^d\right\}
    +O(pt^{d+2})\\
 &=B(v)pt^d\{1+O(p^{d-1})+O(t^2)\},
\end{align*}
where the last line uses $B(v)>0$ fixed. Substitute
$p=\kappa t^{d-4}$ into this formula and \eqref{eq:penaltyjoint}.
Both leading terms then have order $t^{2d-4}$, proving
\eqref{eq:jointlimit}. The optimizing coefficient is obtained by completing
the square; $p(t)\in(0,1/2]$ eventually.

Finally, the residual splitting matrix is
$-\eta\sigma''(b)\E[\chi XX^\top]=0$ by the same moment argument.
The difference of the parity-task risk and the independent-label risk for
any logit $f$ is exactly $-\eta\E[\chi f]$. For a finite sum of smooth
neurons with zero incoming weights, any derivative of this difference in
incoming coordinates of total order less than $d$ is a linear combination
of moments $\E[\chi X^\alpha]$ with $|\alpha|<d$, and vanishes. This proves
(iii). It asserts a derivative identity, not an impossibility for arbitrary
sample-based algorithms, active probing, or optimization away from the parent.
\end{proof}

The decreasing path approaches an embedding with a zero-output child at
$t=0$. It is not a small perturbation of the equal-weight duplication in
the full widened parameter vector: its outgoing proportions differ by order
one from $(1/2,1/2)$. The matched constraints are width, outgoing absolute
sum, incoming radius, and weighted incoming mean, not full parameter
displacement, initialization, or optimization cost.

The theorem compares positive operations at the same width. It does not
say that a vanishing weight is the unique or practical remedy. A fixed
nonzero radius, a different activation, a changed intercept, untied
retraining, signed coefficients, or a different architecture can alter
the conclusion. In particular, the higher-order descent can be extremely
small despite a substantial information gap.

\begin{proposition}[Raw score signal-to-noise cost]\label{prop:snr}
For the same even-parity law, let
$H_t(X)=\sigma(b+tV)-\sigma(b)-t\sigma'(b)V$ and estimate
$m_t=\E[(Y-1/2)H_t(X)]$ by an iid sample average. If $B(v)>0$, then
\[
 m_t=\eta B(v)t^d+O(t^{d+2}),\qquad
 \operatorname{Var}((Y-1/2)H_t)=2K(v)t^4+o(t^4).
\]
Thus the sample size needed for this estimator's standard error to be at
most $m_t$ is asymptotic to
$2K(v)[\eta^2B(v)^2]^{-1}t^{-2d+4}$.
\end{proposition}
\begin{proof}
The mean follows from \eqref{eq:fouriermoment}. Also
$H_t=\sigma''(b)t^2V^2/2+O(t^3)$, while $(Y-1/2)^2=1/4$
identically. Its second moment is therefore
$\sigma''(b)^2\E V^4t^4/16+o(t^4)=2K(v)t^4+o(t^4)$.
Subtracting the squared mean changes only order $t^{2d}$.
For $n$ iid observations the variance of the average is this variance
divided by $n$, which gives the stated standard-error criterion.
\end{proof}
This is an exact asymptotic description of a particular naive score estimator,
not an information-theoretic sample lower bound. If the parity feature is
known, $\E[(Y-1/2)\chi(X)]=\eta$ is estimable without this $t$ penalty.
Knowing or discovering that feature is a different information resource.
